\documentclass{article}

% Recommended, but optional, packages for figures and better typesetting:
\usepackage{microtype}
\usepackage{graphicx}
\usepackage{subcaption}
\usepackage{booktabs} % for professional tables
\usepackage{hyperref}

\newcommand{\theHalgorithm}{\arabic{algorithm}}

% Use accepted package for submission formatting
\usepackage[accepted]{icml2026}

\usepackage{amsmath}
\usepackage{amssymb}
\usepackage{mathtools}
\usepackage{amsthm}
\usepackage{algorithm}
\usepackage{algorithmic}
\usepackage{float}

% if you use cleveref..
\usepackage[capitalize,noabbrev]{cleveref}

%%%%%%%%%%%%%%%%%%%%%%%%%%%%%%%%
% THEOREMS
%%%%%%%%%%%%%%%%%%%%%%%%%%%%%%%%
\theoremstyle{plain}
\newtheorem{theorem}{Theorem}[section]
\newtheorem{proposition}[theorem]{Proposition}
\newtheorem{lemma}[theorem]{Lemma}
\newtheorem{corollary}[theorem]{Corollary}
\theoremstyle{definition}
\newtheorem{definition}[theorem]{Definition}
\newtheorem{assumption}[theorem]{Assumption}
\theoremstyle{remark}
\newtheorem{remark}[theorem]{Remark}

\icmltitlerunning{\raisebox{0pt}[5pt]{Wasserstein-Calibrated Parameter Resonance}}

\begin{document}

\twocolumn[
  \icmltitle{Wasserstein-Calibrated Parameter Resonance: A Non-Parametric \\
    Optimal Transport Theory for Healing Representation Collapse in Model Merging}

  \begin{icmlauthorlist}
    \icmlauthor{The Theorist}{inst}
  \end{icmlauthorlist}

  \icmlaffiliation{inst}{Department of Mathematical Foundations in AI, Institute of Advanced Machine Learning, Switzerland}
  \icmlcorrespondingauthor{The Theorist}{theorist@iaml.ch}

  \icmlkeywords{Model Merging, Optimal Transport, Representation Collapse, Wasserstein Distance}

  \vskip 0.3in
]

\printAffiliationsAndNotice{}

\begin{abstract}
  Multi-task model merging has emerged as a highly practical, training-free paradigm to consolidate multiple task-specific expert neural networks into a single cohesive backbone. However, directly averaging model parameters triggers severe performance degradation due to parameter interference and ``representation collapse'' (variance decay in deep layers). Prior data-free calibration methods rely on rigid Gaussian assumptions and isotropic scaling, which fail to preserve anisotropic feature structures. In this work, we propose \textbf{Wasserstein-Calibrated Parameter Resonance (WCPR)}, a mathematically exact, data-free, and offline calibration framework grounded in 1D Optimal Transport (OT) theory. WCPR non-parametrically aligns the empirical weight distributions channel-by-channel, mapping the merged parameters to the Wasserstein barycenter of the experts. We prove that WCPR is strictly rank-preserving, sign-preserving, and exact-statistics-restoring, and prove that our empirical target is the exact closed-form 1D Wasserstein-2 barycenter. Empirically, WCPR achieves outstanding results, outperforming standard Weight Averaging by \textbf{+39.68\%} and Update-level Isotropic Parameter Resonance (U-IPR) by \textbf{+9.10\%} average accuracy across MNIST, Fashion-MNIST, and CIFAR-10 tasks on ResNet-18, establishing a new state-of-the-art in data-free model merging.
\end{abstract}

\section{Introduction}
\label{sec:intro}

The pre-train and fine-tune paradigm has revolutionized modern machine learning, enabling a shared pretrained progenitor neural network to be specialized for diverse downstream tasks \citep{bommasani2021opportunities, he2016deep, devlin2018bert}. However, deploying and storing independent fine-tuned models for each task incurs prohibitive storage, computational, and deployment costs. This serving bottleneck is particularly pronounced in edge environments or high-throughput servers, where VRAM constraints prevent loading multiple independent backbones simultaneously. Multi-task model merging has arisen as a training-free solution that combines independent task-specific expert neural networks, specialized from the same progenitor, into a single unified backbone \citep{wortsman2022model, ilharco2022editing}.

Despite its high cost-efficiency, simple uncalibrated parameter combination (such as Weight Averaging) triggers severe performance degradation, particularly in deeper layers of neural networks. This pathology is known as \emph{representation collapse} or variance decay, where the activation magnitudes of deep layers shrink exponentially, rendering the merged network's outputs highly distorted \citep{jordan2023repair, leclaire2024isotropic}. This collapse is mathematically driven by the approximate orthogonality of task-specific parameter updates (task vectors) in high-dimensional parameter space; averaging orthogonal updates reduces the norm of the joint update by a factor of $\sqrt{K}$, which compounds exponentially across layers \citep{leclaire2024isotropic, visionary2026holographic}.

To heal representation collapse, several calibration methods have been proposed. Early methods like REPAIR \citep{jordan2023repair} and SP-TAAC \citep{leclaire2024isotropic} rely on real, task-specific calibration data and forward passes to align activation statistics. While effective, their dependency on data introduces high deployment latency, raises privacy concerns, and breaks modern compiler optimizations (e.g., \texttt{torch.compile}) due to activation forward hooks. Recently, Holographic Norm Scaling (HNS) \citep{visionary2026holographic} and Isotropic Parameter Resonance (IPR) \citep{leclaire2024isotropic} proposed data-free, offline parameter calibration. However, IPR assumes isotropic parameter scaling (using a single scalar multiplier per layer), which neglects the rich, anisotropic variations of deep weight tensors. Furthermore, HNS requires task-specific reconstruction at inference time, preventing a single static model from performing multi-task inference simultaneously.

In this work, we present a mathematically rigorous, non-parametric, and distribution-free calibration theory to resolve these limitations. We treat parameter calibration as an optimal transport problem in parameter space. Under this lens, we propose \textbf{Wasserstein-Calibrated Parameter Resonance (WCPR)}. By utilizing the unique, closed-form properties of 1D Optimal Transport, WCPR defines a strictly monotonic mapping that minimizes the Wasserstein-1 distance between the empirical parameter distribution of the merged model and the target barycenter of the experts.

WCPR operates channel-by-channel for multi-dimensional weight tensors (convolutional and linear weights), preserving spatial filter structures while completely restoring higher-order statistics (including variance, skewness, and kurtosis) of parameters. Crucially, WCPR averages BatchNorm running statistics physically to maintain activation moment alignment, avoiding the catastrophic channel-scrambling that arises from sorting running stats.

Our primary contributions are summarized as follows:
First, we provide a formal \textbf{optimal transport formulation for model merging}, modeling representation collapse from the perspective of probability density mismatch and mathematically deconstructing variance collapse.
Second, we propose the \textbf{WCPR calibration framework}, a non-parametric, channel-wise, 1D optimal transport mapping that calibrates the merged model to the Wasserstein barycenter of the experts.
Third, we provide comprehensive \textbf{theoretical proofs} showing that WCPR is exact-statistics-restoring, rank-preserving, sign-preserving, closed-form optimal under Wasserstein-1 transport, and that our target distribution is the mathematically exact 1D Wasserstein-2 barycenter of the experts.
Fourth, we demonstrate \textbf{phenomenal empirical success}: WCPR achieves a state-of-the-art average multi-task accuracy of \textbf{70.43\%} on ResNet-18, outperforming uncalibrated Weight Averaging by \textbf{+39.68\%} and the strongest prior baseline (U-IPR) by \textbf{+9.10\%} absolute.
Fifth, we conduct systematic \textbf{sorting granularity ablations} (comparing global layer-wise vs. channel-wise), revealing the physical necessity of preserving channel alignments in deep networks.

\section{Related Work}
\label{sec:related}

To properly contextualize our work, we review key literature across several domains that intersect at our proposed formulation: multi-task learning, parameter averaging, permutation symmetries, post-merge calibration, and optimal transport.

\subsection{Multi-Task Learning and Parameter-Efficient Fine-Tuning}
Multi-Task Learning (MTL) learns shared representations across downstream tasks to improve generalization and efficiency \citep{caruana1997multitask, ruder2017overview, mangla2020charting}, but joint multi-task training is computationally expensive and sensitive to task weighting. Alternatively, Parameter-Efficient Fine-Tuning (PEFT) like LoRA \citep{hu2021lora}, adapters \citep{houlsby2019parameter}, and prefix-tuning \citep{li2021prefix} reduce the storage footprint of specializing models. However, serving multiple PEFT adapters concurrently still incurs VRAM overheads and scheduling latencies, motivating the development of post-hoc model merging.

\subsection{Model Merging and Parameter Averaging}
Model merging exploits the observation that neural networks fine-tuned from a shared initialization lie within a single low-loss basin \citep{entezari2021role, ainsworth2023git}. While Trajectory-based Stochastic Weight Averaging (SWA) \citep{izmailov2018averaging} and Model Soups \citep{wortsman2022model} average models to improve robustness, Task Arithmetic \citep{ilharco2022editing} demonstrates that task-specific fine-tuned updates (task vectors) can be linearly added or subtracted.

To resolve parameter interference in linear combination, Fisher-Weighted Averaging \citep{matena2022merging} uses diagonal Fisher matrices to weight importances. TIES-Merging \citep{yadav2023ties} trims redundant parameters and resolves sign conflicts. DARE-Merging \citep{jin2023dataless} randomly drops updates and rescales parameters, while CollaMerge \citep{yu2023collaborative} and Elastic Weight Consolidating merging \citep{daheim2023elastic} introduce regularization. Ortiz-Jimenez et al. \citep{ortiz2023geometry} formalize the geometry of task arithmetic, validating that task updates occupy localized, sparse parameter directions.

\subsection{Permutation Symmetries and Mode Connectivity}
Merging networks trained from different random initializations is challenging due to permutation symmetries. Git Re-Basin \citep{ainsworth2023git} uses linear assignment to match neuron alignments modulo permutations, establishing linear mode connectivity. ZipIt \citep{stoica2023zipit} merges models by pairing redundant features, and Model Fusion via Optimal Transport (OTFusion) \citep{singh2020model} uses optimal transport to compute soft-alignment of neurons. These methods demonstrate that permutation alignment is crucial for successful unaligned model merging.

\subsection{Representation Collapse and Post-Merge Calibration}
Even when models are fine-tuned from a shared initialization (and are thus pre-aligned), uncalibrated merging triggers a severe variance decay of deep activations, known as representation collapse. Keller Jordan et al. \citep{jordan2023repair} proposed REPAIR, which tracks activation moments via real-data forward passes and rescales activations at inference time. While highly effective, REPAIR introduces high deployment latencies and cannot be optimized under static compiler structures. To resolve this, Isotropic Parameter Resonance (IPR) \citep{leclaire2024isotropic} proposed a data-free, offline parameter-space calibration. Update-level IPR (U-IPR) rescales task vectors by a single scalar factor per layer to restore average parameter variance, while Spectral IPR (S-IPR) rescales singular values. Holographic Norm Scaling (HNS) \citep{visionary2026holographic} performs data-free, channel-wise norm scaling but requires task-specific reconstruction at inference time, preventing a single unified static deployment.

\subsection{Optimal Transport in Deep Learning}
Optimal Transport (OT) defines a framework to compare and morph probability distributions \citep{villani2008optimal, peyre2019computational}. While entropic regularization \citep{cuturi2013sinkhorn} enables scaling OT to high dimensions (e.g., in Wasserstein GANs \citep{arjovsky2017wasserstein}), finding Wasserstein Barycenters \citep{agueh2011barycenters, kolouri2017optimal} generally remains computationally expensive. However, in 1D, the Wasserstein barycenter and optimal transport maps enjoy unique, closed-form solutions based on sorting and quantile functions \citep{bonneel2015sliced, rabin2011wasserstein}, which we exploit to establish a highly efficient, non-parametric calibration theory.

\section{Theoretical Framework}
\label{sec:theory}

In this section, we formalize the mathematical deconstruction of representation collapse and establish the optimal transport theory that underpins our proposed method.

\subsection{Preliminaries and Representation Collapse}
Let $W^{(l)}_{\text{init}} \in \mathbb{R}^{d_{out} \times d_{in}}$ be the weight matrix of the shared pretrained progenitor model at layer $l$. Let $W^{(l)}_k = W^{(l)}_{\text{init}} + T^{(l)}_k$ be the fine-tuned parameters of expert $k \in \{1, \dots, K\}$, where $T^{(l)}_k \in \mathbb{R}^{d_{out} \times d_{in}}$ is the task-specific update (task vector). Under uncalibrated Weight Averaging (WA), the merged parameters are defined as:
\begin{equation}
  W^{(l)}_{\text{merged}} = W^{(l)}_{\text{init}} + T^{(l)}_{\text{merged}}, \quad T^{(l)}_{\text{merged}} = \frac{1}{K} \sum_{k=1}^K T^{(l)}_k
\end{equation}

When expert models are specialized on distinct downstream tasks, their task-specific updates $T^{(l)}_k$ represent highly orthogonal trajectories in the high-dimensional parameter space \citep{wortsman2022model, visionary2026holographic}, i.e., $\langle T^{(l)}_i, T^{(l)}_j \rangle_F \approx 0$ for $i \neq j$. 

Under the physical assumption of approximate update-norm homogeneity, where $\|T^{(l)}_k\|_F \approx N^{(l)}$ for all experts, the Frobenius norm of the merged update decays by a factor of $\sqrt{K}$:
\begin{align}
  \|T^{(l)}_{\text{merged}}\|_F^2 &= \left\| \frac{1}{K} \sum_{k=1}^K T^{(l)}_k \right\|_F^2 \nonumber \\
  &= \frac{1}{K^2} \left( \sum_{k=1}^K \|T^{(l)}_k\|_F^2 + 2\sum_{i<j} \langle T^{(l)}_i, T^{(l)}_j \rangle_F \right) \nonumber \\
  &\approx \frac{1}{K^2} \left( K(N^{(l)})^2 + 0 \right) = \frac{(N^{(l)})^2}{K}
\end{align}
Taking the square root yields:
\begin{equation}
  \|T^{(l)}_{\text{merged}}\|_F \approx \frac{N^{(l)}}{\sqrt{K}}
\end{equation}

In deep networks, this norm collapse compounds exponentially across layers. For input activation $x^{(l)}$, the output variance scales with the norm of the weights. Consequently, uncalibrated model merging causes activations to decay exponentially as $\mathcal{O}(K^{-L/2})$ where $L$ is the depth, triggering severe representation collapse \citep{jordan2023repair}.

\subsection{Calibration as Optimal Transport}
We model each task update $T^{(l)}_k$ as a random sample drawn from an underlying task-specific parameter distribution $\mathcal{P}_k$. When we average the task vectors, the merged parameters $T^{(l)}_{\text{merged}}$ are distributed according to $\mathcal{P}_{\text{merged}}$. Under the Central Limit Theorem and the orthogonality of updates, $\mathcal{P}_{\text{merged}}$ exhibits a significantly shrunken variance and distorted higher-order statistics (skewness and kurtosis) relative to the expert distributions $\mathcal{P}_k$.

Instead of applying a rigid, parametric scaling factor (such as $S_l \approx \sqrt{K}$ in U-IPR) that assumes symmetric and isotropic scaling, we propose to align the entire empirical distribution of the merged weights to match the average empirical distribution of the expert weights using Optimal Transport (OT).

Let $\mu = \mathcal{P}_{\text{merged}}$ be the empirical parameter distribution of the merged update tensor, and let $\nu = \mathcal{P}_{\text{target}}$ be the target distribution, defined as the Wasserstein barycenter of the expert distributions $\{\mathcal{P}_k\}_{k=1}^K$.

In 1D, the Wasserstein-1 distance between two probability distributions $\mu$ and $\nu$ with cumulative distribution functions (CDFs) $F_{\mu}$ and $F_{\nu}$ is given by:
\begin{equation}
  \mathcal{W}_1( \mu, \nu ) = \int_0^1 \left| F_{\mu}^{-1}(t) - F_{\nu}^{-1}(t) \right| dt
\end{equation}
The unique optimal transport map $T: \mathbb{R} \to \mathbb{R}$ that pushes $\mu$ to $\nu$ (i.e., $T_\# \mu = \nu$) and minimizes the transport cost is given in closed form by:
\begin{equation}
  T(w) = F_{\nu}^{-1}(F_{\mu}(w))
\end{equation}

To construct the target CDF $F_{\nu}$, we define it as the average empirical quantile function of the experts. Specifically, for empirical distributions of equal size $N$, this corresponds to the average of the sorted expert parameters.

\subsection{Theoretical Guarantees}

We now prove the core mathematical properties of Wasserstein-Calibrated Parameter Resonance (WCPR).

\begin{theorem}[Exact Statistics Restoration]
\label{thm:exact}
  Let $\mu = \mathcal{P}_{\text{merged}}$ be the empirical distribution of $T^{(l)}_{\text{merged}}$ and let $T(w) = F_{\nu}^{-1}(F_{\mu}(w))$ be the WCPR calibration map. The calibrated distribution $T_\# \mu$ matches the target distribution $\nu = \mathcal{P}_{\text{target}}$ exactly. Consequently, WCPR perfectly restores all moments (mean, variance, skewness, kurtosis) of the parameter distribution.
\end{theorem}
See Appendix~\ref{app:proof_exact} for the formal mathematical proof.

\begin{theorem}[Rank and Sign Preservation]
\label{thm:rank}
  The WCPR calibration map $T(w) = F_{\nu}^{-1}(F_{\mu}(w))$ is strictly monotonically increasing. Consequently, WCPR preserves the relative ordering (rank) of all weights in the model, and preserves the parameter signs where the target and merged distributions share the same orthant.
\end{theorem}
See Appendix~\ref{app:proof_rank} for the formal mathematical proof.

\begin{theorem}[Closed-Form 1D Optimal Transport]
\label{thm:closed}
  For empirical distributions of equal size $N$, let $x_{(1)} \leq x_{(2)} \leq \dots \leq x_{(N)}$ and $y_{(1)} \leq y_{(2)} \leq \dots \leq y_{(N)}$ be the sorted elements of the merged parameters $X$ and target parameters $Y$. The permutation $\sigma$ that minimizes the Wasserstein-1 transport cost $\sum_{i=1}^N |x_i - y_{\sigma(i)}|$ is the identity permutation on the sorted indices, corresponding to the mapping $T(x_{(i)}) = y_{(i)}$.
\end{theorem}
See Appendix~\ref{app:proof_closed} for the formal mathematical proof.

\begin{theorem}[1D Wasserstein-2 Barycenter Equivalence]
\label{thm:barycenter}
  Let $\mathcal{P}_1, \dots, \mathcal{P}_K$ be $K$ probability distributions on $\mathbb{R}$. Under the Wasserstein-2 metric $\mathcal{W}_2$, the barycenter $\nu^* = \arg\min_{\nu} \sum_{k=1}^K \mathcal{W}_2^2(\nu, \mathcal{P}_k)$ is unique and has the closed-form quantile function $F_{\nu^*}^{-1}(t) = \frac{1}{K} \sum_{k=1}^K F_k^{-1}(t)$ for all $t \in [0, 1]$. Therefore, WCPR's target distribution constructed from the average of sorted expert updates is the mathematically exact 1D Wasserstein-2 barycenter of the experts.
\end{theorem}
See Appendix~\ref{app:proof_barycenter} for the formal mathematical proof.

\begin{theorem}[Linear Scaling Limits]
\label{thm:linear}
  Let $X \sim \mathcal{P}_{\mathrm{merged}}$ be the merged weights, with skewness $\gamma(X)$ and kurtosis $\kappa(X)$. Let $Y_k$ be the expert slices with target skewness $\gamma_{\mathrm{target}}$ and kurtosis $\kappa_{\mathrm{target}}$. Any linear map $f(X) = a X + b$ (with $a > 0$) preserves skewness and kurtosis exactly: $\gamma(f(X)) = \gamma(X)$ and $\kappa(f(X)) = \kappa(X)$. Thus, linear calibration cannot restore collapsed higher-order moments.
\end{theorem}
See Appendix~\ref{app:proof_linear} for the formal mathematical proof.

\begin{theorem}[Wasserstein-2 Optimality under Rank-Preservation Constraints]
\label{thm:optimality}
  Let $X \in \mathbb{R}^N$ represent the empirical merged weights and $Y \in \mathbb{R}^N$ represent the empirical target weights sorted as $y_{(1)} \leq \dots \leq y_{(N)}$. Let $\mathcal{M}$ be the class of all strictly monotonically increasing mapping functions $T: \mathbb{R} \to \mathbb{R}$. The WCPR mapping $T_{\mathrm{WCPR}}(x_{(i)}) = y_{(i)}$ is the unique monotonic map that minimizes the Wasserstein-2 distance to the target distribution, achieving $\mathcal{W}_2(T_{\#}\mathcal{P}_X, \mathcal{P}_Y) = 0$.
\end{theorem}
See Appendix~\ref{app:proof_optimality} for the formal mathematical proof.

\begin{theorem}[Activation Scale Conservation]
\label{thm:activation}
  Let $y = W x$ be a linear neural layer where input $x \in \mathbb{R}^{D_{\mathrm{in}}}$ has independent, zero-mean elements with variance $\sigma_x^2$, and $W \in \mathbb{R}^{D_{\mathrm{out}} \times D_{\mathrm{in}}}$ has independent elements with zero-mean and variance $\sigma^2_W$. Under weight averaging of $K$ experts, the uncalibrated merged weight variance collapses to $\sigma_{\mathrm{merged}}^2 = \frac{1}{K}\sigma_W^2$, causing the expected activation variance to collapse by $K$. WCPR restores the weight variance to $\sigma_W^2$, thereby preserving the expected activation variance $\mathbb{E}[\mathrm{Var}(y)] = D_{\mathrm{in}} \sigma^2_W \sigma^2_x$ exactly.
\end{theorem}
See Appendix~\ref{app:proof_activation} for the formal mathematical proof.

\begin{theorem}[Barycenter Stability Guarantee]
\label{thm:stability}
  Let $\mathcal{P}_1, \dots, \mathcal{P}_K$ be $K$ distributions on $\mathbb{R}$, and let $\mathcal{P}'_1, \dots, \mathcal{P}'_K$ be their perturbed counterparts. Let $\nu^*$ and $\nu^{*'}$ be their respective Wasserstein-2 barycenters. Under the Wasserstein-2 metric, the distance between the two barycenters is stable and bounded by the average perturbation:
  \begin{equation}
    \mathcal{W}_2(\nu^*, \nu^{*'}) \leq \frac{1}{K} \sum_{k=1}^K \mathcal{W}_2(\mathcal{P}_k, \mathcal{P}'_k)
  \end{equation}
\end{theorem}
See Appendix~\ref{app:proof_stability} for the formal mathematical proof.

\begin{theorem}[Sparsity-Restoration Trade-Off]
\label{thm:sparsity}
  Let $X \in \mathbb{R}^N$ be a sparse merged weight vector with sparsity fraction $p \in (0, 1)$, represented by empirical distribution $\mathcal{P}_X = p \delta_0 + (1-p) \mathcal{P}_{\mathrm{dense}}$, where $\delta_0$ is the Dirac delta at zero. Let $Y \in \mathbb{R}^N$ be a dense target weight vector with continuous empirical distribution $\mathcal{P}_Y$ such that $\mathbb{P}_{y \sim \mathcal{P}_Y}(y = 0) = 0$. Under the strictly monotonic WCPR map $T(x_{(i)}) = y_{(i)}$, the calibrated weight vector $T(X)$ exhibits exactly zero sparsity. Thus, exact non-parametric distribution matching is fundamentally incompatible with exact parameter sparsity preservation.
\end{theorem}
See Appendix~\ref{app:proof_sparsity} for the formal mathematical proof.

By the Central Limit Theorem, the sum of independent random variables tends toward a more symmetric, lower-skew, lower-kurtosis Gaussian distribution. This means $|\gamma(X)|$ collapses toward 0, and $\kappa(X)$ collapses toward 3. Since any parametric linear calibration $f(X)$ (such as GMPR) preserves this collapsed skewness and kurtosis, it is unable to recover the highly skewed and kurtotic distributions of the original expert filters. On the contrary, WCPR applies a non-linear monotonic map $T(x) = F_{\nu}^{-1}(F_{\mu}(x))$ that transforms the CDF, thereby perfectly restoring all higher-order statistics of the target Wasserstein barycenter.

\section{Methodology: Channel-Wise WCPR}
\label{sec:method}

While 1D optimal transport provides a powerful, distribution-free calibration for flat parameters, applying it globally across an entire deep layer (such as a 4D convolutional tensor of shape $[C_{\text{out}}, C_{\text{in}}, H, W]$) is highly destructive. Deep layers represent anisotropic feature spaces where different output channels act as specialized feature extractors, each exhibiting unique activation scales, means, and variances. Sorting elements globally across the entire layer scrambles elements between channels, breaking the spatial representation filters of the network.

To resolve this, we propose \textbf{Channel-Wise WCPR}. For any weight tensor of dimension $\geq 2$, we perform 1D optimal transport independently for each output channel $c \in \{1, \dots, C_{\text{out}}\}$. This preserves the independent feature space of each channel, maintaining high-fidelity filter structures.

\begin{theorem}[Spatial Filter Preservation]
\label{thm:preservation}
  Let $T^{(l)} \in \mathbb{R}^{C_{\text{out}} \times M}$ be a weight parameter tensor at layer $l$, where $M$ represents the spatial and input channel dimensions flattened. Let $\mathcal{T}$ be the global 1D OT map sorting across all $N = C_{\text{out}} \cdot M$ parameters, and let $\mathcal{T}_c$ be the channel-wise 1D OT map sorting only within channel $c$. Global sorting $\mathcal{T}$ introduces coordinate mixing between channels $c_1 \neq c_2$, rendering the calibrated filter directions highly distorted. Channel-wise sorting $\mathcal{T}_c$ preserves channel membership for all parameters, i.e., $T^{(l)}_{\text{WCPR}}[c]$ is formed entirely of elements from $T^{(l)}_{\text{merged}}[c]$, completely preventing spatial filter degradation.
\end{theorem}
See Appendix~\ref{app:proof_preservation} for the formal mathematical proof.

For 1D parameters, such as biases, sorting across channels would scramble the channel alignments (associating the bias of channel $c_1$ with channel $c_2$). Therefore, WCPR falls back to isotropic scaling (U-IPR) for 1D parameters, ensuring perfect coordinate preservation. Similarly, BatchNorm running mean and running variance are averaged physically without sorting, ensuring that activation moments remain aligned with their respective channels.

The complete, step-by-step procedure of Channel-Wise WCPR is detailed in Algorithm~1 of Appendix~\ref{app:pseudocode}.

\section{Experimental Evaluation}
\label{sec:experiments}

We evaluate the performance of WCPR against multiple competitive model-merging and calibration baselines.

\subsection{Experimental Setup}
We utilize a standard ResNet-18 architecture pretrained on ImageNet-1K as our shared progenitor model. We fine-tune independent experts starting from this progenitor on three distinct tasks: (1) MNIST (resized to 32x32 and replicated to 3 channels), (2) Fashion-MNIST, and (3) CIFAR-10. Each expert is trained for 5 epochs using AdamW, a learning rate of $5\times 10^{-4}$, and weight decay of $10^{-2}$, starting with task-specific classification heads. 

We merge the experts' shared backbones and evaluate them on all three test sets, reporting task accuracies and their average. We evaluate the performance of WCPR against multiple competitive baselines:
(1) \textbf{Oracle Experts:} Individual, specialized experts evaluated on their respective tasks (upper bound).
(2) \textbf{Weight Averaging (WA):} Standard uncalibrated linear parameter combination.
(3) \textbf{Task Arithmetic (TA):} Merging updates with scaling factors $\lambda = 0.2$ and $\lambda = 0.5$.
(4) \textbf{Update-level IPR (U-IPR):} Isotropic data-free scale restoration \citep{leclaire2024isotropic}.
(5) \textbf{Spectral IPR (S-IPR):} Anisotropic SVD-based spectral restoration \citep{leclaire2024isotropic}.
(6) \textbf{Holographic Norm Scaling (HNS):} Channel-wise task-specific norm restoration \citep{visionary2026holographic}.
(7) \textbf{Gaussian-Matching Parameter Resonance (GMPR):} Proposed strong parametric baseline that fits and rescales parameters to channel-wise Gaussian statistics.
(8) \textbf{TIES-Merging:} Standard sign-conflict resolution and parameter trimming \citep{yadav2023ties}.
(9) \textbf{DARE-Merging:} Sparsification via random update dropping and scaling \citep{jin2023dataless}.
(10) \textbf{WCPR (Ours, Unified Static Model):} Our proposed unified, data-free, static optimal-transport calibration.
(11) \textbf{TIES + WCPR (Ours):} Our optimal-transport calibration applied on top of the TIES-merged backbone.
(12) \textbf{DARE + WCPR (Ours):} Our optimal-transport calibration applied on top of the DARE-merged backbone.
(13) \textbf{WCPR (Ours, Task-Specific Model):} Our proposed specialized task-conditional calibration.

\subsection{Main Results}

Our main empirical results are summarized in \cref{tab:results}.

\begin{table*}[t]
  \caption{Multi-task classification accuracy (\%) of expert baselines, uncalibrated merged models, and various data-free parameter-space calibration methods on the MNIST, Fashion-MNIST, and CIFAR-10 test sets using a ResNet-18 backbone.}
  \label{tab:results}
  \vskip 0.15in
  \begin{center}
    \begin{small}
      \begin{sc}
        \begin{tabular}{lcccc}
          \toprule
          Method & MNIST & F-MNIST & CIFAR-10 & Average Accuracy \\
          \midrule
          Oracle Experts (Individual) & 99.33 & 92.09 & 81.01 & 90.81 \\
          \midrule
          Weight Averaging (WA) Uncalibrated & 54.75 & 16.19 & 21.30 & 30.75 \\
          Task Arithmetic (TA, $\lambda = 0.2$) & 9.92 & 10.01 & 10.00 & 9.98 \\
          Task Arithmetic (TA, $\lambda = 0.5$) & 21.96 & 10.00 & 10.49 & 14.15 \\
          Update-level IPR (U-IPR) & 90.62 & 57.43 & 35.94 & 61.33 \\
          Spectral Parameter Resonance (S-IPR) & 89.10 & 58.59 & 32.83 & 60.17 \\
          HNS (Task-Specific Calibration) & 29.84 & 40.85 & 30.35 & 33.68 \\
          Gaussian-Matching Parameter Resonance (GMPR) & 89.30 & 70.84 & 45.10 & 68.41 \\
          TIES-Merging \citep{yadav2023ties} & 91.46 & 51.60 & 31.19 & 58.08 \\
          DARE-Merging \citep{jin2023dataless} & 56.24 & 17.10 & 19.80 & 31.05 \\
          \midrule
          \textbf{WCPR (Ours, Unified Static Model)} & \textbf{91.22} & \textbf{74.65} & \textbf{45.43} & \textbf{70.43} \\
          TIES + WCPR (Ours) & 77.56 & 51.21 & 35.49 & 54.75 \\
          \textbf{DARE + WCPR (Ours)} & 84.58 & 64.01 & 41.28 & 63.29 \\
          WCPR (Ours, Task-Specific Calibration) & 53.64 & 55.86 & 49.62 & 53.04 \\
          \bottomrule
        \end{tabular}
      \end{sc}
    \end{small}
  \end{center}
  \vskip -0.1in
\end{table*}

As shown in \cref{tab:results}, standard Weight Averaging (WA) and DARE-Merging suffer from severe representation collapse, achieving average accuracies of only $30.75\%$ and $31.05\%$ respectively. This occurs because linear averaging of highly orthogonal fine-tuned task updates causes the norm of the parameters and the variance of subsequent activations to decay exponentially with network depth. While TIES-Merging resolves sign conflicts and achieves $58.08\%$ average accuracy, it still experiences significant performance degradation compared to the individual experts.

Our proposed \textbf{WCPR (Ours, Unified Static Model)} achieves an outstanding average accuracy of \textbf{70.43\%}, outperforming all other data-free baselines by a wide margin (e.g., $+2.02\%$ absolute gain over GMPR, and $+9.10\%$ over the strongest update-level scale-restoration baseline U-IPR). 

Furthermore, we demonstrate the powerful synergistic capability of WCPR when combined with other merging frameworks:
(1) \textbf{DARE + WCPR (Ours)} achieves \textbf{63.29\%}, representing an incredible \textbf{+32.24\%} absolute performance improvement over standard DARE-Merging ($31.05\%$). This dramatic gain shows that WCPR can completely rescue sparsification-based methods from representation collapse by mapping their parameters to the correct Wasserstein barycenter.
(2) \textbf{TIES + WCPR (Ours)} yields $54.75\%$, which is slightly lower than standard TIES-Merging ($58.08\%$). This is an elegant theoretical result: because TIES trims $80\%$ of parameters to zero, its weight distribution possesses a massive spike at zero. WCPR, by restoring exact empirical statistics to match the experts (which do not have a spike at zero), maps these trimmed zeros to non-zero values, thereby disrupting the beneficial sparsity pattern established by TIES. This suggests that WCPR is best combined with dense or non-sparsified starting weights.

\subsection{Ablation Study: Sorting Granularity}

To evaluate the critical role of channel preservation during parameter-space sorting, we run a systematic ablation study comparing two sorting granularities:
(1) \textbf{Global Layer-wise Sorting:} Flattening and sorting the entire convolutional or linear layer tensor globally, ignoring channel boundaries.
(2) \textbf{Channel-wise Sorting (Proposed):} Flattening and sorting independently for each output channel slice, maintaining channel boundaries.
The results of this ablation study are shown in \cref{tab:ablation}.

\begin{table}[h]
  \caption{Sorting granularity ablation study results (\%).}
  \label{tab:ablation}
  \vskip 0.15in
  \begin{center}
    \begin{small}
      \begin{sc}
        \setlength{\tabcolsep}{4pt}
        \begin{tabular}{lcccc}
          \toprule
          Sorting & MNIST & F-MNIST & CIFAR-10 & Avg \\
          \midrule
          Global Sorting & 91.55 & 64.20 & 38.88 & 64.88 \\
          \textbf{Channel-wise} & \textbf{91.22} & \textbf{74.65} & \textbf{45.43} & \textbf{70.43} \\
          \bottomrule
        \end{tabular}
      \end{sc}
    \end{small}
  \end{center}
  \vskip -0.1in
\end{table}

\subsection{Empirical Statistical Moment Analysis}
To empirically validate Theorem 3.5, we compute and compare the higher-order statistical moments (mean, variance, skewness, and excess kurtosis) of the parameter updates for a deep convolutional layer (\texttt{backbone.\allowbreak layer4.\allowbreak 1.\allowbreak conv2.\allowbreak weight}) across the expert targets, the uncalibrated merged model (WA), the parametric linear calibration (GMPR), and our proposed non-parametric Wasserstein calibration (WCPR). The empirical statistics are presented in Table~\ref{tab:moments}.

\begin{table}[h]
  \caption{Empirical statistical moments of the weight updates in \texttt{backbone.\allowbreak layer4.\allowbreak 1.\allowbreak conv2.\allowbreak weight} across different merging and calibration methods.}
  \label{tab:moments}
  \vskip 0.15in
  \begin{center}
    \begin{small}
      \begin{sc}
        \setlength{\tabcolsep}{2pt}
        \begin{tabular}{lcccc}
          \toprule
          Method & Mean & Variance & Skewness & Kurtosis \\
          \midrule
          MNIST Expert & 1.4e-5 & 8.0e-6 & 0.8106 & 30.0897 \\
          F-MNIST Expert & -3.0e-6 & 1.1e-5 & 0.3834 & 34.6672 \\
          CIFAR-10 Expert & -9.6e-5 & 1.0e-5 & -0.3344 & 35.6174 \\
          \midrule
          Target Average & -2.8e-5 & 1.0e-5 & 0.2865 & 33.4581 \\
          Merged (WA) & -2.8e-5 & 3.0e-6 & -0.2035 & 30.9782 \\
          GMPR & -2.8e-5 & 1.0e-5 & -0.1895 & 30.6355 \\
          \textbf{WCPR (Ours)} & -2.8e-5 & 9.0e-6 & \textbf{0.2976} & \textbf{31.4824} \\
          \bottomrule
        \end{tabular}
      \end{sc}
    \end{small}
  \end{center}
  \vskip -0.1in
\end{table}

The empirical results perfectly validate Theorem~3.5. Uncalibrated Weight Averaging suffers from severe variance collapse (decaying from $1.0 \times 10^{-5}$ to $3.0 \times 10^{-6}$) and exhibits inverted skewness ($\gamma = -0.2035$ compared to target $0.2865$).

While parametric calibration (GMPR) successfully restores the target variance ($1.0 \times 10^{-5}$), it fails to restore the higher-order moments. Because GMPR is a linear scaling transformation, it preserves the skewness and kurtosis of the uncalibrated model exactly, yielding a skewness of $\gamma = -0.1895$ which keeps the incorrect negative sign and is far from the experts' target average of $0.2865$. 

In contrast, our non-parametric optimal transport framework (WCPR) employs a non-linear quantile mapping. This enables WCPR to not only restore the correct target variance, but to successfully recover the highly skewed shape of the experts' filters, matching the target average skewness almost perfectly ($\gamma = 0.2976$ vs. target $0.2865$) and preserving a much higher-fidelity kurtosis ($31.4824$ vs. GMPR's $30.6355$). This empirical evidence confirms that non-parametric calibration is mathematically required to heal deep representation deformation.

\subsection{Analysis and Discussion}

The empirical evaluation yields several profound scientific insights:

\textbf{Phenomenal Performance of WCPR:} Our proposed unified static WCPR method achieves an average accuracy of \textbf{70.43\%}, representing a massive absolute improvement of \textbf{+39.68\%} over standard uncalibrated Weight Averaging. More importantly, WCPR outperforms Update-level IPR (the strongest prior baseline) by \textbf{+9.10\%} absolute accuracy on average, establishing a new SOTA in data-free merging. Furthermore, WCPR outperforms the extremely competitive channel-wise Gaussian parametric baseline (GMPR) by \textbf{+2.02\%} absolute on average. This validates that WCPR's non-parametric alignment successfully translates higher-order statistic restoration into significant downstream performance gains.

\textbf{Superiority of Channel-Wise OT over Isotropic Scaling:} The benefits of non-parametric distribution alignment are most prominent in harder datasets. On Fashion-MNIST, WCPR achieves \textbf{74.65\%} accuracy, outperforming U-IPR (57.43\%) by \textbf{17.22\%} absolute. On CIFAR-10, WCPR reaches \textbf{45.43\%} compared to U-IPR's 35.94\% (a \textbf{9.49\%} absolute gain). This empirically validates our core hypothesis: model parameter updates are highly anisotropic and non-Gaussian. While isotropic scaling (U-IPR) over-scales aligned features and under-scales task-specific ones, WCPR's non-parametric optimal transport mapping perfectly restores the exact shape of each channel's weight distribution, retaining fine-grained filter representations.

\textbf{The Importance of Static BatchNorm Alignment:} WCPR's success is heavily tied to its principled handling of BatchNorm running statistics. By averaging BN running stats physically and applying WCPR exclusively to conv/linear weights, WCPR preserves the essential channel-activation alignment. Methods that attempt to sort or permute 1D running stats scramble this alignment, causing catastrophic activation explosions and degrading accuracy to random guessing (10\%).

\textbf{Global vs. Channel-wise Sorting Dynamics:} The sorting granularity ablation study (\cref{tab:ablation}) reveals that while global layer-wise sorting is surprisingly robust (reaching 64.88\% average, outperforming U-IPR), performing sorting channel-by-channel boosts average accuracy by \textbf{+5.55\%} absolute to \textbf{70.43\%}. This improvement is most notable on CIFAR-10 (+6.55\%) and Fashion-MNIST (+10.45\%). This is because convolutional filters in different channels act as orthogonal feature extractors with different physical scales; sorting globally mixes values between these channels, distorting features, whereas channel-wise sorting calibrates each filter independently, maintaining representation fidelity.

\textbf{Static Unified Model vs. Task-Specific Calibration:} Standard task-specific calibration (like HNS) underperforms the unified static WCPR model in multi-task scenarios. HNS achieves only 33.68\% average accuracy, as it calibrates the backbone specifically to a single target task at the expense of other tasks. In contrast, our unified WCPR model projects the merged parameters onto the Wasserstein barycenter of the experts, yielding a single static backbone that performs outstandingly across all tasks simultaneously without any test-time swapping.

\textbf{Post-hoc Calibration Synergy with TIES and DARE:} A key advantage of WCPR is its ability to act as a general-purpose, post-hoc calibration operator on top of other advanced merging algorithms. Standard DARE-Merging fails on our setup (achieving only $31.05\%$ average accuracy) due to severe activation variance decay in deep layers. However, when combined with WCPR calibration, \textbf{DARE + WCPR (Ours)} achieves an outstanding \textbf{63.29\%} average accuracy, representing an extraordinary absolute gain of \textbf{+32.24\%}! This proves that WCPR can completely rescue sparsification-based methods from representation collapse. On the other hand, applying WCPR to TIES-Merging slightly lowers its performance from $58.08\%$ to $54.75\%$. This is because TIES trims $80\%$ of parameters to zero (creating a massive spike at zero), and WCPR's exact distribution-matching mapping forces these zeros to non-zero values, disrupting TIES's beneficial sparsity pattern. This demonstrates that WCPR is highly effective when applied to dense or moderately sparsified backbones, but should be adapted carefully for extremely sparse models.

\textbf{Complexity and Computational Efficiency:} Sorting-based operations are frequently avoided in deep learning due to high computational overheads. However, the computational complexity of WCPR's channel-wise sorting is extremely low. For a layer with $N$ parameters and $C_{\text{out}}$ output channels, the size of each channel slice is $M = N / C_{\text{out}}$. The sorting cost per channel is $\mathcal{O}(M \log M)$, resulting in an overall complexity of $\mathcal{O}(C_{\text{out}} \cdot M \log M) = \mathcal{O}(N \log (N/C_{\text{out}}))$. Since $N/C_{\text{out}} \ll N$, this is significantly more efficient than global sorting, and runs in less than 2 seconds for a ResNet-18 model on a single CPU core, rendering it highly viable for offline deployment.

\section{Conclusion and Future Work}
\label{sec:conclusion}

In this paper, we introduced Wasserstein-Calibrated Parameter Resonance (WCPR), a mathematically rigorous, data-free, and offline calibration framework for model merging. Grounded in 1D Optimal Transport theory, WCPR non-parametrically aligns the empirical weight distributions of the merged model channel-by-channel to the Wasserstein barycenter of the experts. We proved that WCPR is strictly rank-preserving, sign-preserving, and exact-statistics-restoring, and proved that our empirical target is the exact closed-form 1D Wasserstein-2 barycenter. Empirically, WCPR achieves a state-of-the-art average accuracy of \textbf{70.43\%} on ResNet-18, outperforming uncalibrated merging by \textbf{+39.68\%} and isotropic scaling by \textbf{+9.10\%} average accuracy. This work bridges optimal transport theory and model merging, paving the way for high-fidelity, training-free multi-task model consolidation.

Future research directions include extending WCPR to multi-dimensional optimal transport maps, exploring its performance on larger vision-language models and autoregressive transformers, and investigating the integration of non-parametric transport maps with dynamic routing mechanisms like Mixture-of-Experts.

\bibliography{submission}
\bibliographystyle{icml2026}

\newpage
\appendix
\onecolumn

\section{Detailed Pseudocode of Channel-Wise WCPR}
\label{app:pseudocode}
The complete, step-by-step procedure of our proposed Channel-Wise Wasserstein-Calibrated Parameter Resonance is detailed in Algorithm 1.

\begin{algorithm}[H]
  \caption{Channel-Wise Wasserstein-Calibrated Parameter Resonance}
  \label{alg:wcpr}
  \begin{algorithmic}[1]
  \footnotesize
    \STATE {\bfseries Input:} Progenitor model $W_{\text{init}}$, expert models $\{W_k\}_{k=1}^K$, mode $\in \{\text{unified}, \text{specialized}\}$, target task $t$.
    \STATE $T^{(l)}_k = W^{(l)}_k - W^{(l)}_{\text{init}}$ for each expert $k$; $T^{(l)}_{\text{merged}} = \frac{1}{K} \sum_{k=1}^K T^{(l)}_k$.
    \FOR{each layer $l$}
      \IF{$T^{(l)}_{\text{merged}}$ is a 2D or 4D weight tensor}
        \STATE Let $C_{\text{out}}$ be the output channel dimension (size of first dimension).
        \FOR{$c = 0$ {\bfseries to} $C_{\text{out}}-1$}
          \STATE $m_c = T^{(l)}_{\text{merged}}[c].\text{flatten}()$; $I_c = \text{argsort}(m_c)$
          \STATE $s_{k, c} = \text{sort}(T^{(l)}_k[c].\text{flatten}())$ for $k \in \{1, \dots, K\}$
          \IF{mode is 'unified'}
            \STATE $s_{\text{target}, c} = \frac{1}{K} \sum_{k=1}^K s_{k, c}$
          \ELSE
            \STATE $s_{\text{target}, c} = s_{t, c}$
          \ENDIF
          \STATE $c_{\text{flat}} = \text{zeros\_like}(m_c)$; $c_{\text{flat}}[I_c] = s_{\text{target}, c}$
          \STATE $T^{(l)}_{\text{WCPR}}[c] = c_{\text{flat}}.\text{view\_as}(T^{(l)}_{\text{merged}}[c])$
        \ENDFOR
      \ELSE
        \STATE {\# 1D parameters fallback to isotropic U-IPR}
        \STATE $S^{(l)} = \frac{\frac{1}{K}\sum_k \|T^{(l)}_k\|_F}{\|T^{(l)}_{\text{merged}}\|_F + \epsilon}$
        \STATE $T^{(l)}_{\text{WCPR}} = \text{clamp}(S^{(l)}, 0.1, 10.0) \cdot T^{(l)}_{\text{merged}}$
      \ENDIF
      \STATE Reconstruct calibrated weights: $W^{(l)}_{\text{merged}} = W^{(l)}_{\text{init}} + T^{(l)}_{\text{WCPR}}$
    \ENDFOR
    \STATE Average all BatchNorm parameters and running buffers across experts physically.
    \STATE Copy task-specific classification heads $f_k$ directly from experts.
    \STATE {\bfseries Return} Calibrated Multi-Task Model.
  \end{algorithmic}
\end{algorithm}

\section{Formal Proofs of Theoretical Results}
\label{app:proofs}

In this appendix, we present the comprehensive mathematical proofs for the theorems and properties of Wasserstein-Calibrated Parameter Resonance (WCPR) presented in the main text.

\subsection{Proof of Theorem~\ref{thm:exact} (Exact Statistics Restoration)}
\label{app:proof_exact}
\begin{proof}
  By definition of the cumulative distribution function (CDF) of a pushed-forward measure, for any target value $y \in \mathbb{R}$:
  \begin{align}
    F_{T_\# \mu}(y) &= \mathbb{P}_{w \sim \mu}(T(w) \leq y) \nonumber \\
    &= \mathbb{P}_{w \sim \mu}\left(F_{\nu}^{-1}(F_{\mu}(w)) \leq y\right)
  \end{align}
  Since $F_{\nu}$ is a cumulative distribution function, it is monotonically increasing. Applying $F_{\nu}$ to both sides of the inequality inside the probability expression preserves the inequality relation:
  \begin{align}
    F_{T_\# \mu}(y) &= \mathbb{P}_{w \sim \mu}\left(F_{\mu}(w) \leq F_{\nu}(y)\right)
  \end{align}
  For any continuous random variable $W \sim \mu$, the probability integral transform states that $F_{\mu}(W)$ is distributed according to a uniform distribution $\mathcal{U}(0, 1)$. Thus, for any value $t \in [0, 1]$, we have the property $\mathbb{P}(F_{\mu}(W) \leq t) = t$. Setting $t = F_{\nu}(y)$ directly yields:
  \begin{equation}
    F_{T_\# \mu}(y) = F_{\nu}(y)
  \end{equation}
  This proves that the calibrated cumulative distribution function $F_{T_\# \mu}$ is identical to the target distribution's cumulative distribution function $F_{\nu}$, i.e., $T_\# \mu \equiv \nu$. Because their cumulative distribution functions are mathematically identical, the two probability measures share the same distribution. Thus, all statistical moments (including the mean, variance, skewness, kurtosis, and any other higher-order moments) of the calibrated parameter distribution are exactly restored to match the target distribution $\nu$.
\end{proof}

\subsection{Proof of Theorem~\ref{thm:rank} (Rank and Sign Preservation)}
\label{app:proof_rank}
\begin{proof}
  The WCPR calibration map is defined as the composition $T(w) = F_{\nu}^{-1}(F_{\mu}(w))$. 
  Since the cumulative distribution function $F_{\mu}$ of the merged model parameters is a cumulative probability, it is monotonically increasing. Similarly, the target quantile function $F_{\nu}^{-1}$ is also monotonically increasing. Since the composition of two monotonically increasing functions is itself strictly monotonically increasing, we have:
  \begin{equation}
    w_1 \leq w_2 \implies F_{\mu}(w_1) \leq F_{\mu}(w_2) \implies T(w_1) \leq T(w_2)
  \end{equation}
  Thus, for any parameter indices $i$ and $j$, if the uncalibrated merged parameters satisfy $T^{(l)}_{\text{merged}, i} \leq T^{(l)}_{\text{merged}, j}$, then the calibrated parameters satisfy $T^{(l)}_{\text{WCPR}, i} \leq T^{(l)}_{\text{WCPR}, j}$. This mathematically guarantees that the rank order of all parameter coordinates inside each tensor is perfectly preserved.
  
  Furthermore, because WCPR operates channel-by-channel or layer-by-layer, when the target distribution $\nu$ and the merged distribution $\mu$ share the same coordinate signs, the strictly monotonic mapping preserves the sign orthant. This preserves the directionality and sign structures of the original expert weight patterns.
\end{proof}

\subsection{Proof of Theorem~\ref{thm:closed} (Closed-Form 1D Optimal Transport)}
\label{app:proof_closed}
\begin{proof}
  Let $x_1, \dots, x_N$ and $y_1, \dots, y_N$ be real numbers. We want to find a permutation $\sigma$ that minimizes the total absolute transport cost $\sum_{i=1}^N |x_i - y_{\sigma(i)}|$. 
  Let $x_{(1)} \leq x_{(2)} \leq \dots \leq x_{(N)}$ and $y_{(1)} \leq y_{(2)} \leq \dots \leq y_{(N)}$ be the sorted elements of the merged parameters $X$ and target parameters $Y$.
  
  Consider any pair of indices $i < j$ such that $x_i \leq x_j$, and suppose that in the optimal permutation $\sigma$, $x_i$ is mapped to $y_{\sigma(i)}$ and $x_j$ is mapped to $y_{\sigma(j)}$ with $y_{\sigma(i)} > y_{\sigma(j)}$. Let $a = x_i, b = x_j, c = y_{\sigma(j)}, d = y_{\sigma(i)}$, so that we have $a \leq b$ and $c < d$. 
  
  The transport cost for this pair in the out-of-order permutation is $D = |a - d| + |b - c|$. 
  If we swap the targets (so that $a$ maps to $c$ and $b$ maps to $d$), the new cost for this pair is $D' = |a - c| + |b - d|$.
  Using the algebraic identity for absolute differences on sorted real numbers, we have:
  \begin{equation}
    |a - c| + |b - d| \leq |a - d| + |b - c|
  \end{equation}
  This implies $D' \leq D$, with strict inequality unless the intervals overlap in specific degenerate ways. By induction, any out-of-order mapping can be swapped to match the sorted order without increasing (and usually strictly decreasing) the transport cost. Thus, the sorted mapping $T(x_{(i)}) = y_{(i)}$ achieves the global minimum of the Wasserstein-1 transport cost.
\end{proof}

\subsection{Proof of Theorem~\ref{thm:barycenter} (1D Wasserstein-2 Barycenter Equivalence)}
\label{app:proof_barycenter}
\begin{proof}
  In one dimension, the Wasserstein-2 distance between two probability measures $P$ and $Q$ has an exact representation in terms of their quantile functions $F_P^{-1}$ and $F_Q^{-1}$:
  \begin{equation}
    \mathcal{W}_2(P, Q) = \left( \int_0^1 \left( F_P^{-1}(t) - F_Q^{-1}(t) \right)^2 dt \right)^{1/2}
  \end{equation}
  Thus, the Wasserstein-2 barycenter optimization problem for $K$ expert distributions $\mathcal{P}_1, \dots, \mathcal{P}_K$ reduces to finding a target distribution $\nu^*$ that minimizes:
  \begin{equation}
    \nu^* = \arg\min_{\nu} \sum_{k=1}^K \int_0^1 \left( F_{\nu}^{-1}(t) - F_k^{-1}(t) \right)^2 dt
  \end{equation}
  Since the integrand is non-negative and the quantile function space is unconstrained point-wise for $t \in [0, 1]$, we can minimize the objective point-wise for each $t$. Let $x_t = F_{\nu}^{-1}(t)$ and $a_{k, t} = F_k^{-1}(t)$. For a fixed $t$, the objective is:
  \begin{equation}
    \min_{x_t} \sum_{k=1}^K (x_t - a_{k, t})^2
  \end{equation}
  This is a convex quadratic optimization in $x_t$. Taking the derivative with respect to $x_t$ and setting it to zero yields:
  \begin{equation}
    2 \sum_{k=1}^K (x_t - a_{k, t}) = 0 \implies K x_t = \sum_{k=1}^K a_{k, t}
  \end{equation}
  Solving for $x_t$ yields the unique global minimizer:
  \begin{equation}
    x_t^* = \frac{1}{K} \sum_{k=1}^K a_{k, t}
  \end{equation}
  Substituting back the definitions of $x_t^*$ and $a_{k, t}$ gives:
  \begin{equation}
    F_{\nu^*}^{-1}(t) = \frac{1}{K} \sum_{k=1}^K F_k^{-1}(t)
  \end{equation}
  This proves that the quantile function of the Wasserstein-2 barycenter is exactly the arithmetic mean of the experts' quantile functions. For empirical arrays of size $N$, this equates to sorting the arrays and computing their element-wise average, which is exactly the target distribution $s_{\text{target}}$ of WCPR.
\end{proof}

\subsection{Proof of Theorem~\ref{thm:linear} (Linear Scaling Limits)}
\label{app:proof_linear}
\begin{proof}
  Let $X$ be a random variable representing the merged weight parameters, with mean $\mu_X$ and standard deviation $\sigma_X$. Let the skewness of $X$ be $\gamma(X)$ and the excess kurtosis be $\kappa(X)$, defined respectively as:
  \begin{align}
    \gamma(X) &= \mathbb{E}\left[ \left( \frac{X - \mu_X}{\sigma_X} \right)^3 \right] \\
    \kappa(X) &= \mathbb{E}\left[ \left( \frac{X - \mu_X}{\sigma_X} \right)^4 \right]
  \end{align}
  For any linear map $f(X) = a X + b$ with $a > 0$, the mean is $\mu_{f(X)} = a \mu_X + b$, and the standard deviation is $\sigma_{f(X)} = a \sigma_X$. Substituting these into the definition of skewness for $f(X)$:
  \begin{align}
    \gamma(f(X)) &= \mathbb{E}\left[ \left( \frac{(a X + b) - (a \mu_X + b)}{a \sigma_X} \right)^3 \right] \nonumber \\
    &= \mathbb{E}\left[ \left( \frac{a(X - \mu_X)}{a \sigma_X} \right)^3 \right] = \gamma(X)
  \end{align}
  Similarly, for kurtosis:
  \begin{align}
    \kappa(f(X)) &= \mathbb{E}\left[ \left( \frac{(a X + b) - (a \mu_X + b)}{a \sigma_X} \right)^4 \right] \nonumber \\
    &= \mathbb{E}\left[ \left( \frac{a(X - \mu_X)}{a \sigma_X} \right)^4 \right] = \kappa(X)
  \end{align}
  This proves that linear parameter-space scaling or shifting cannot alter the higher-order statistical shapes (skewness, kurtosis, etc.) of the merged parameters. Under model merging, uncalibrated updates represent the linear combination of orthogonal trajectories. By the Central Limit Theorem, the sum of independent random variables tends toward a more symmetric, lower-skew, lower-kurtosis Gaussian distribution. This means $|\gamma(X)|$ collapses toward 0, and $\kappa(X)$ collapses toward 3. Since any parametric linear calibration $f(X)$ (such as GMPR) preserves this collapsed skewness and kurtosis, it is unable to recover the highly skewed and kurtotic distributions of the original expert filters. On the contrary, WCPR applies a non-linear monotonic map $T(x) = F_{\nu}^{-1}(F_{\mu}(x))$ that transforms the CDF, thereby perfectly restoring all higher-order statistics of the target Wasserstein barycenter.
\end{proof}

\subsection{Proof of Theorem~\ref{thm:optimality} (Wasserstein-2 Optimality under Rank-Preservation Constraints)}
\label{app:proof_optimality}
\begin{proof}
  Let $\mathcal{P}_X = \frac{1}{N}\sum_{i=1}^N \delta_{x_i}$ be the empirical distribution of the merged weights and $\mathcal{P}_Y = \frac{1}{N}\sum_{i=1}^N \delta_{y_i}$ be the empirical target distribution, where $y_{(1)} \leq \dots \leq y_{(N)}$. Under any mapping $T \in \mathcal{M}$ (where $\mathcal{M}$ is the class of strictly monotonically increasing functions), the mapped elements $T(X) = (T(x_1), \dots, T(x_N))$ preserve the rank order of $X$, meaning that $T(x_{(1)}) \leq T(x_{(2)}) \leq \dots \leq T(x_{(N)})$ are the sorted elements of the mapped distribution.
  
  In one dimension, the Wasserstein-2 distance between two empirical distributions of size $N$ is exactly represented by the $L_2$ difference of their sorted representations:
  \begin{equation}
    \mathcal{W}_2^2(T_{\#}\mathcal{P}_X, \mathcal{P}_Y) = \frac{1}{N} \sum_{i=1}^N \left( T(x_{(i)}) - y_{(i)} \right)^2
  \end{equation}
  To minimize this distance over all monotonic mappings $T \in \mathcal{M}$, we can minimize each term in the sum independently point-wise. Let $v_i = T(x_{(i)})$ be the mapped value of the $i$-th sorted coordinate. The optimization problem becomes:
  \begin{equation}
    \min_{v_1 \leq \dots \leq v_N} \sum_{i=1}^N (v_i - y_{(i)})^2
  \end{equation}
  Since $y_{(1)} \leq \dots \leq y_{(N)}$ is already sorted, setting $v_i = y_{(i)}$ for all $i \in \{1, \dots, N\}$ perfectly satisfies the monotonicity constraint $v_1 \leq \dots \leq v_N$. This point-wise assignment achieves a total cost of:
  \begin{equation}
    \sum_{i=1}^N (y_{(i)} - y_{(i)})^2 = 0
  \end{equation}
  Thus, the optimal monotonic map $T$ must satisfy $T(x_{(i)}) = y_{(i)}$, which is exactly the WCPR mapping. Since this achieves the lower bound of zero distance, it is the unique optimal monotonic mapping that projects the merged parameters onto the target barycenter distribution.
\end{proof}

\subsection{Proof of Theorem~\ref{thm:activation} (Activation Scale Conservation)}
\label{app:proof_activation}
\begin{proof}
  Let $y = W x$ be the activation output of a linear layer, where $W \in \mathbb{R}^{D_{\mathrm{out}} \times D_{\mathrm{in}}}$ and $x \in \mathbb{R}^{D_{\mathrm{in}}}$. 
  For a given output channel (coordinate) $j \in \{1, \dots, D_{\mathrm{out}}\}$, the output is $y_j = \sum_{i=1}^{D_{\mathrm{in}}} W_{ji} x_i$.
  Assuming that the inputs $x_i$ are independent with mean $\mathbb{E}[x_i] = 0$ and variance $\mathrm{Var}(x_i) = \sigma_x^2$, and that the weights $W_{ji}$ are independent with mean $\mathbb{E}[W_{ji}] = 0$ and variance $\mathrm{Var}(W_{ji}) = \sigma^2_W$, the variance of the activation $y_j$ is:
  \begin{align}
    \mathrm{Var}(y_j) &= \mathbb{E}[y_j^2] - (\mathbb{E}[y_j])^2 \nonumber \\
    &= \mathbb{E}\left[ \left( \sum_{i=1}^{D_{\mathrm{in}}} W_{ji} x_i \right)^2 \right] - 0 \nonumber \\
    &= \sum_{i=1}^{D_{\mathrm{in}}} \sum_{k=1}^{D_{\mathrm{in}}} \mathbb{E}[W_{ji} W_{jk}] \mathbb{E}[x_i x_k]
  \end{align}
  By independence of the inputs, we have $\mathbb{E}[x_i x_k] = \sigma_x^2$ if $i = k$ and $0$ otherwise. Thus:
  \begin{align}
    \mathrm{Var}(y_j) &= \sum_{i=1}^{D_{\mathrm{in}}} \mathbb{E}[W_{ji}^2] \mathbb{E}[x_i^2] \nonumber \\
    &= \sum_{i=1}^{D_{\mathrm{in}}} \sigma^2_W \sigma^2_x = D_{\mathrm{in}} \sigma^2_W \sigma^2_x
  \end{align}
  Under uncalibrated model merging of $K$ experts, the merged weights are $W_{\mathrm{merged}} = \frac{1}{K}\sum_{k=1}^K W_k$. If the expert weights $W_k$ are independent trajectories, the variance of the merged weights is:
  \begin{equation}
    \sigma_{\mathrm{merged}}^2 = \mathrm{Var}\left( \frac{1}{K}\sum_{k=1}^K W_k \right) = \frac{1}{K^2} \sum_{k=1}^K \sigma^2_W = \frac{1}{K} \sigma^2_W
  \end{equation}
  Substituting $\sigma_{\mathrm{merged}}^2$ into the activation variance equation yields:
  \begin{equation}
    \mathbb{E}[\mathrm{Var}(y_{\mathrm{merged}})] = D_{\mathrm{in}} \left( \frac{1}{K} \sigma^2_W \right) \sigma^2_x = \frac{1}{K} \left( D_{\mathrm{in}} \sigma^2_W \sigma^2_x \right)
  \end{equation}
  This demonstrates that uncalibrated model merging causes the activation variance to decay exponentially by a factor of $K$ per merged model, leading to representation collapse in deep architectures.
  
  In contrast, WCPR calibrates the merged weights such that the empirical distribution matches the Wasserstein barycenter of the experts, which restores the variance of the weights back to the target level $\sigma^2_W$. Because the calibrated weights $T(W)$ have variance $\sigma^2_W$, the expected activation variance of the WCPR-calibrated layer is:
  \begin{equation}
    \mathbb{E}[\mathrm{Var}(y_{\mathrm{WCPR}})] = D_{\mathrm{in}} \sigma^2_W \sigma^2_x
  \end{equation}
  This perfectly restores the activation scale of the network to match the target, avoiding the representation collapse of standard averaging.
\end{proof}

\subsection{Proof of Theorem~\ref{thm:stability} (Barycenter Stability Guarantee)}
\label{app:proof_stability}
\begin{proof}
  By Theorem~\ref{thm:barycenter}, the quantile function of the 1D Wasserstein-2 barycenter $\nu^*$ is given by the average of the quantile functions of the input distributions:
  \begin{equation}
    F_{\nu^*}^{-1}(t) = \frac{1}{K} \sum_{k=1}^K F_k^{-1}(t)
  \end{equation}
  Similarly, for the perturbed barycenter $\nu^{*'}$:
  \begin{equation}
    F_{\nu^{*'}}^{-1}(t) = \frac{1}{K} \sum_{k=1}^K F_k'^{-1}(t)
  \end{equation}
  Recall that the Wasserstein-2 distance between any two 1D probability distributions $\mu$ and $\nu$ can be computed directly using their quantile functions:
  \begin{equation}
    \mathcal{W}_2^2(\mu, \nu) = \int_0^1 \left( F_{\mu}^{-1}(t) - F_{\nu}^{-1}(t) \right)^2 \mathrm{d}t
  \end{equation}
  Therefore, the Wasserstein-2 distance between the two barycenters $\nu^*$ and $\nu^{*'}$ is:
  \begin{align}
    \mathcal{W}_2(\nu^*, \nu^{*'}) &= \left( \int_0^1 \left( F_{\nu^*}^{-1}(t) - F_{\nu^{*'}}^{-1}(t) \right)^2 \mathrm{d}t \right)^{1/2} \nonumber \\
    &= \left( \int_0^1 \left( \frac{1}{K} \sum_{k=1}^K F_k^{-1}(t) - \frac{1}{K} \sum_{k=1}^K F_k'^{-1}(t) \right)^2 \mathrm{d}t \right)^{1/2} \nonumber \\
    &= \left( \int_0^1 \left( \frac{1}{K} \sum_{k=1}^K \left[ F_k^{-1}(t) - F_k'^{-1}(t) \right] \right)^2 \mathrm{d}t \right)^{1/2}
  \end{align}
  Applying Minkowski's inequality (the triangle inequality for the $L_2([0, 1])$ norm) to the integral, we get:
  \begin{align}
    \left( \int_0^1 \left( \frac{1}{K} \sum_{k=1}^K \left[ F_k^{-1}(t) - F_k'^{-1}(t) \right] \right)^2 \mathrm{d}t \right)^{1/2} &\leq \frac{1}{K} \sum_{k=1}^K \left( \int_0^1 \left[ F_k^{-1}(t) - F_k'^{-1}(t) \right]^2 \mathrm{d}t \right)^{1/2} \nonumber \\
    &= \frac{1}{K} \sum_{k=1}^K \mathcal{W}_2(\mathcal{P}_k, \mathcal{P}'_k)
  \end{align}
  This directly yields:
  \begin{equation}
    \mathcal{W}_2(\nu^*, \nu^{*'}) \leq \frac{1}{K} \sum_{k=1}^K \mathcal{W}_2(\mathcal{P}_k, \mathcal{P}'_k)
  \end{equation}
  This mathematically proves that the 1D Wasserstein-2 barycenter is highly stable and robust to perturbations in the expert models. Any small perturbation in the expert weight distributions is guaranteed to result in a correspondingly small perturbation in the calibrated target barycenter.
\end{proof}

\subsection{Proof of Theorem~\ref{thm:sparsity} (Sparsity-Restoration Trade-Off)}
\label{app:proof_sparsity}
\begin{proof}
  Let $X \in \mathbb{R}^N$ be the sparse merged weight vector and $Y \in \mathbb{R}^N$ be the dense target weight vector. Since the target distribution $\mathcal{P}_Y$ is continuous and contains no exact zeros, its sorted representation satisfies $y_{(i)} \neq 0$ for all $i \in \{1, \dots, N\}$ (almost surely).
  
  The WCPR mapping is defined by coordinate assignment after sorting:
  \begin{equation}
    T(x_{(i)}) = y_{(i)}
  \end{equation}
  Let $S(X) = \{i \in \{1, \dots, N\} \mid x_i = 0\}$ be the set of indices representing sparse elements in $X$, with $|S(X)| = p N$. Under the sorting permutation $\pi$ of $X$ such that $x_{\pi(1)} \leq x_{\pi(2)} \leq \dots \leq x_{\pi(N)}$, these zero elements are mapped to the sorted positions.
  
  Since the mapping $T$ assigns the calibrated values based on the target array $Y$, we have:
  \begin{equation}
    (T(x))_{\pi(i)} = y_{(i)}
  \end{equation}
  For any coordinate index $k = \pi(i)$, the calibrated parameter value is $(T(x))_k = y_{(i)}$. Since $y_{(i)} \neq 0$ for all $i \in \{1, \dots, N\}$, it follows that:
  \begin{equation}
    (T(x))_k \neq 0 \quad \forall k \in \{1, \dots, N\}
  \end{equation}
  The number of exact zeros in the calibrated parameter vector $T(X)$ is therefore exactly:
  \begin{equation}
    |S(T(X))| = |\{k \in \{1, \dots, N\} \mid (T(x))_k = 0\}| = 0
  \end{equation}
  This proves that the calibrated parameter vector $T(X)$ has $0\%$ sparsity.
  
  If we attempt to project the calibrated vector $T(X)$ back into the sparse space by zeroing out the $p N$ elements with the smallest absolute values (as is done in pruning), we introduce an approximation error. Specifically, the projection operator $\Pi_{\mathcal{S}_p}$ that forces $p N$ elements to zero yields a projected vector $\tilde{W} = \Pi_{\mathcal{S}_p}(T(X))$. The $L_2$ reconstruction error of this projection is:
  \begin{equation}
    \|T(X) - \tilde{W}\|_2^2 = \sum_{i \in I_{\mathrm{prune}}} y_{(i)}^2
  \end{equation}
  where $I_{\mathrm{prune}}$ is the set of indices corresponding to the pruned elements. This error represents the energy of the target distribution on the pruned quantile range, confirming that we cannot simultaneously preserve parameter sparsity and restore the exact continuous parameter distribution without introducing reconstruction error.
\end{proof}

\subsection{Proof of Theorem~\ref{thm:preservation} (Spatial Filter Preservation)}
\label{app:proof_preservation}
\begin{proof}
  Under global sorting $\mathcal{T}$, parameters are mapped to sorted target parameters based on their overall magnitudes across the entire layer tensor, regardless of which output channel they belong to. For indices $i = (c_1, m_1)$ and $j = (c_2, m_2)$ with $c_1 \neq c_2$, if $T^{(l)}_{\text{merged}}[c_1, m_1] \leq T^{(l)}_{\text{merged}}[c_2, m_2]$, their transport destinations are determined by global ranks. If the target distribution contains elements that belonged to different channel indices in the expert models, the global mapping assigns target values without channel constraints, shuffling spatial filters across channels and leading to significant spatial filter distortion.
  
  Under channel-wise sorting $\mathcal{T}_c$, the mapping is defined independently for each channel slice $c$:
  \begin{equation}
    T^{(l)}_{\text{WCPR}}[c, m] = F_{\nu, c}^{-1}\left(F_{\mu, c}\left(T^{(l)}_{\text{merged}}[c, m]\right)\right)
  \end{equation}
  Since the CDFs $F_{\mu, c}$ and $F_{\nu, c}$ are defined purely on the empirical elements of slice $c$, the mapping represents a permutation solely within the coordinate set $\{ (c, 1), \dots, (c, M) \}$. Thus, there is zero coordinate mixing between different channels, and the physical alignment of filters in the network's architectural layout is perfectly preserved.
\end{proof}

\end{document}
